\documentclass[12pt]{article}

\usepackage{amsmath}
\usepackage[margin=1in]{geometry}
\usepackage{fancyhdr}
\usepackage{amsthm}
\usepackage{amssymb}
\usepackage{pgfplots}
\usepackage{tikz}
\usepackage{bbm}
\usepackage{bm}
\usetikzlibrary{shapes.geometric}
\usepackage{enumerate}
\usepackage[shortlabels]{enumitem}

\pgfplotsset{compat=1.18}

\newcommand{\pderiv}[2][]{\frac{\partial#1}{\partial#2}}

\pagestyle{fancy}
\fancyhead[l]{Hudson Benites}
\fancyhead[c]{Day 8}
\fancyfoot[c]{\thepage}
\renewcommand{\headrulewidth}{0.2pt}
\setlength{\headheight}{15pt}
\fancyhead[r]{\today}

\newcommand{\magnitude}[1]{\left\lvert #1 \right\rvert} %chktex 1

\begin{document}
\begin{enumerate}[label=\Alph*., leftmargin=*]
    \item \begin{enumerate}
        \item Consider \vspace{2in}
        \item We can define $\alpha:\mathbb{R}\to\mathbb{R}^{3}$ as\begin{equation*}
            \alpha(t)= (f (t),f (-t),42)
        \end{equation*} where\begin{equation*}
            f(t)=\begin{cases}
                0&\text{if }t\leq 0\\
                e^{-t^{-1}}&\text{if }t>0
            \end{cases}.
        \end{equation*}
    \end{enumerate}
    \item \begin{enumerate}
        \item The issue with Gary's proof is that it is not guaranteed that $\gamma$ is on $S$. $\gamma$ must be on $S$ on a small interval $(-\epsilon, \epsilon)\subseteq \mathbb{R}$, but $\gamma$ is only guaranteed to be on $S$ when $t=0$.
        \item \begin{proof}
            Let $X\in T_{p}S$. Then, we have an $\alpha: (a,b)\to\mathbb{R}^n$ in $S$ such that $\alpha(0)=p$ and $\dot{\alpha}(0)=X$. Consider $\tilde{\alpha}:(\frac{a}{c},\frac{b}{c})\to \mathbb{R}^{n}$ given by $\tilde{\alpha}(t)=\alpha(ct)$, $c\neq0$. We see that ${\tilde{\alpha}}(0)=p$. We compute that $\dot{\tilde{\alpha}}(t)=c\cdot\dot{\alpha}(ct)$. So, when $t=0$,
            $\dot{\tilde{\alpha}}(0)=c\cdot\dot{\alpha}(0)=cX$. We now show that we also have $\mathbf{0}$. We know we have $\tilde{\alpha}(t)=\alpha(-t)$. So, we have $-X\in T_{p}S$. We are given that $T_{p}S$ is closed under addition, so $X-X=0\in T_{p}S$. We are done. 
        \end{proof}
    \end{enumerate}
    \item \begin{proof}
        When $\mu=(0,0,1)$, $k\mu\cdot x=k\cos u$. So, $f(x)=ce^{k\cos u}$. We compute\begin{equation*}
            c\iint_{\mathbb{S}^2}e^{k\cos u}\sin u\, \mathrm{d}u \, \mathrm{d}v=c\int_{-\pi}^{\pi} \int_{0}^{\pi} e^{k\cos u}\sin u \, \mathrm{d}u \, \mathrm{d}v=\frac{2c\pi(e^k-\frac{1}{e^k})}{k}.
        \end{equation*} using a u-substitution where $u=\cos u$. In order to make this distribution physically valid, we must solve for $c$ such that \begin{equation*}
            \frac{2c\pi(e^k-\frac{1}{e^k})}{k}=1.
        \end{equation*} Note that $e^k-\frac{1}{e^k}=e^k-e^{-k}$ which is equivalent to $\sinh{k}$. So, we find that $c=\frac{k}{4\pi\sinh k}$.
    \end{proof}
\end{enumerate}
\end{document}